%%%%%%%%%%%%%%%%%%%%%%%%%%%%%%%%%%%%%%%%%%%%%%%%%%%%%%%%%%%%%%%%%%%%%%%%%%%%%%%%
% - Compile: PDFLaTeX => bibtex => PDFLaTeX => PDFLaTeX
%%%%%%%%%%%%%%%%%%%%%%%%%%%%%%%%%%%%%%%%%%%%%%%%%%%%%%%%%%%%%%%%%%%%%%%%%%%%%%%%
\documentclass[%
 reprint,
superscriptaddress,
% groupedaddress,
%unsortedaddress,
%runinaddress,
%frontmatterverbose, 
%preprint,
%preprintnumbers,
%nofootinbib,
%nobibnotes,
%bibnotes,
 amsmath,amssymb,
 aps,
 prl,
%pra,
%prb,
%rmp,
%prstab,
%prstper,
floatfix,
]{revtex4-2}

\usepackage{mathptmx} % Use Times New Roman font
\newcommand{\rev}[1]{{\leavevmode\color{red}#1}}
\usepackage{graphicx}% Include figure files
\usepackage{dcolumn}% Align table columns on decimal point
\usepackage{bm}% bold math

% minimal working environment. Provide example-image
\usepackage{mwe} 

\usepackage{upgreek}
\usepackage{physics}
\usepackage{braket}
% Customization
\usepackage{xcolor}

\usepackage{hyperref}
\hypersetup{colorlinks,
    bookmarks,
	linkcolor={blue!75!black!80!yellow},
	citecolor={blue!75!black!80!yellow},
	urlcolor={blue!75!black!80!yellow}
}


\begin{document}

\preprint{APS/123-QED}

\title{Full Title Here}

\affiliation{Affiliation 1}
\affiliation{Affiliation 2}
\affiliation{Affiliation 3}

\author{First LastA}
\affiliation{Affiliation 1}
\author{First LastB}
\affiliation{Affiliation 1}
\author{First LastC}
\affiliation{Affiliation 1}
\affiliation{Affiliation 2}
\author{First LastD}
\email{firstlastD@nju.edu.cn}
\affiliation{Affiliation 3}
\author{First LastE}
\email{firstlaste@fudan.edu.cn}
\affiliation{Affiliation 2}
\author{First LastF}
\email{firstlastf@psu.edu}
\affiliation{Affiliation 1}

\begin{abstract}

Abstract here. 
\lipsum[1]



\end{abstract}

\pdfbookmark[1]{Title}{title} 

\maketitle

% \section{Introduction}
\emph{Introduction}---Physical systems in the real world are inevitably open in various senses, making non-Hermitian (NH) physics more ubiquitous than its Hermitian counterpart. 
A hallmark of NH physics is the emergence of complex-valued eigenvalues and non-orthogonal eigenstates, which are absent in Hermitian systems  \cite{Kunst2018BiorthogonalBulkBoundaryCorrespondence, El-Ganainy2018NonHermitianPhysicsPT, Kawabata2019SymmetryTopologyNonHermitian, Ashida2020NonHermitianPhysics, Bergholtz2021ExceptionalTopologyNonHermitian, Ding2022NonHermitianTopologyExceptionalpoint}.
\lipsum[2]



\lipsum[3-6]


%%%%%%%%%%%%%%%%%%%%%%%%%%%%%%%%%%%%%%%%%%%%%%%%%%%%%%%%%%%%%%%%%%%%%%%%%%%%%%%%
\begin{figure}[!htb]
\centering
\includegraphics[width=0.49\textwidth]{example-image}
\caption{
    (a) Sample 1. 
    (b) Sample 2
}
\label{fig:Green_fun}
\end{figure}
%%%%%%%%%%%%%%%%%%%%%%%%%%%%%%%%%%%%%%%%%%%%%%%%%%%%%%%%%%%%%%%%%%%%%%%%%%%%%%%%


\emph{Subtitle 1}---\lipsum[5-6]


When $G(\omega)$ acts on the right eigenstate $\ket{\psi_n^\mathrm{R}}$ from the right or on the left eigenstate $\bra{\psi_n^\mathrm{L}}$ from the left, it yields 
\begin{equation}
G(\omega) \ket{\psi_n^\mathrm{R}}    
=  
\frac{1}{\omega-E_n} 
 \ket{\psi_n^\mathrm{R}}    
 \qc 
 \bra{\psi_n^\mathrm{L}}  G(\omega)
=  
\frac{1}{\omega-E_n} 
 \bra{\psi_n^\mathrm{L}}
 .
\label{eq:green_spec}
\end{equation}
Equation~(\ref{eq:green_spec}) demonstrates that the left and right eigenstates of $G(\omega)$ are identical to those of $H$, and the eigenvalues of $G(\omega)$ are directly related to the Hamiltonian's eigenvalues $E_n$ by ${1}/({\omega-E_n})$.

\lipsum[8-10]

\emph{Nonreciprocal NH lattices}---\lipsum[11]



\emph{Discussions and conclusions}---\lipsum[12]




\begin{acknowledgments}
F. L. thanks the support of startup funds from Penn State University and NSF awards 1234567 and 666666.
\end{acknowledgments}


% bibtex_jiaxin: Jiaxin's customized bibliography file
% bibtex: Other bibliography
\bibliography{bibtex_jiaxin, bibtex}


\clearpage 
\appendix


%%%%%%%%%%%%%%%%%%%%%%%%%%%%%%%%%%%%%%%%%%%%%%%%%%%%%%%%%%%%%%%%%%%%%%%%%%%%%%%%
% End Matter
%%%%%%%%%%%%%%%%%%%%%%%%%%%%%%%%%%%%%%%%%%%%%%%%%%%%%%%%%%%%%%%%%%%%%%%%%%%%%%%%
\onecolumngrid
\section*{End Matter}\label{sec:endmatter}

\twocolumngrid

\emph{Appendix A: Title here}---\lipsum[15]

%%%%%%%%%%%%%%%%%%%%%%%%%%%%%%%%%%%%%%%%%%%%%%%%%%%%%%%%%%%%%%%%%%%%%%%%%%%%%%%%
\begin{figure*}[t!]
    \centering
    \includegraphics[width=0.7\textwidth]{example-image-golden}
    \caption{
        Caption here.
    }
    \label{fig:kaisun:7x7}
\end{figure*}
%%%%%%%%%%%%%%%%%%%%%%%%%%%%%%%%%%%%%%%%%%%%%%%%%%%%%%%%%%%%%%%%%%%%%%%%%%%%%%%%


\emph{Appendix B: Title here}---\lipsum[16-25]

\clearpage

% \clearpage 
% 从这里开始修改编号和前缀
% \setcounter{figure}{0} % 重置编号
% \renewcommand{\thefigure}{\arabic{figure}}
% \renewcommand{\figurename}{Extended Data Fig.}

% 强制让超链接系统更新引用目标
% \makeatletter
% \renewcommand{\theHfigure}{extended.\thefigure} % 修改内部超链接的标识
% \makeatother

% \onecolumngrid
\end{document}
